\documentclass[10pt,a4paper]{amsart}
\usepackage[margin=19mm]{geometry}
\usepackage{amsmath,amssymb,mathtools}
\usepackage[scr=rsfs,cal=euler]{mathalfa}
\usepackage{xfrac}
\usepackage[unicode,hidelinks,bookmarks=false]{hyperref}
\newcommand{\ie}{i.e.\ }
\newcommand{\cf}{cf.\ }
\newcommand{\resp}{respectively\ }
\newcommand{\Subsection}{\subsection}
\providecommand{\nobreakdash}{\nobreak\hbox{-}}
\setlength{\emergencystretch}{2em}
\multlinegap0pt
\usepackage{amssymb}
\newtheorem{lemma}{Lemma}
\title{The Image of the Gassner Representation of the Pure Braid Subgroup has Pairwise Free Generators}
\author{G. Makenzie Cosgrove}
\date{Source: arXiv:2204.12469v4 (CC BY 4.0)}
\begin{document}
\maketitle

\noindent\emph{Source and licence.} The Image of the Gassner Representation of the Pure Braid Subgroup has Pairwise Free Generators. Version arXiv:2204.12469v4 (CC BY 4.0), distributed by arXiv under the Creative Commons Attribution 4.0 International licence, http://creativecommons.org/licenses/by/4.0/, as recorded in the arXiv licence metadata for this exact version and shown on the abstract page https://arxiv.org/abs/2204.12469v4. Full licence notice: \texttt{licenses/arxiv-2204.12469-CC-BY-4.0.txt}.\par\smallskip

\noindent\textit{Editorial scope.} Counted unit: the article's lemma lem:cb-xpl (source0.tex lines 129--144). The article's complete original proof of it is delivered with it (source0.tex lines 146--398, one \texttt{proof} environment of 253 lines). No mathematics is written, repaired or re-derived: the statement and the proof are the article's own lines, in the article's own notation, and no retained line is substituted or re-flowed.  No further source-local context is needed: every cross-reference of every retained line resolves inside the two delivered spans.  Nothing was omitted from any retained span, so the package carries no omission record.  No retained line contains a citation, so the wrapper carries no bibliography and no bibliography entry.  No retained line contains a figure environment, an image-inclusion command or drawing code, so no image is delivered and the package declares no asset dependency.  Editorial formatting only, in addition: an AMS \texttt{amsart} preamble that restates the article's own declarations and macros used by the retained lines, namely the environments \texttt{lemma} on separate counters, together with the standard packages those lines need.  The retained environments print in this document as Lemma 1 in order of appearance, whereas the article prints the same items with the numbering of its own full document.  The complete native source of the article as distributed in the arXiv e-print for this exact version is bundled unchanged as \texttt{sources/125015/source0.tex}.
\begin{lemma}\label{lem:cb-xpl}
	Let $n\ge 2$ and
	\begin{equation*}
	A_{i,j},\ 1\le i<j\le n
	\end{equation*}
	be a pure braid generator for some $i,j$. The matrix
	\begin{equation*}
	cb(A_{i,j})\in GL_n(L_n)  
	\end{equation*}
	is equal to
	\begin{eqnarray*}
		I_n &+\quad c_{i-1}((-t_it_j+t_i)_{i\to j}) &+\quad c_i((t_j-1)_{i\to j})\\
		&+\quad c_{j-1}((t_it_j-t_j)_{i\to j}) &+\quad c_{j}((-t_i+1)_{i\to j}),
	\end{eqnarray*}
	with the following caveat: When $i=1$ the term $c_{i-1}$ is absent.
\end{lemma}

\begin{proof}
	When $n=2$, we need only check $A_{1,2}$, which we obtain as the result of a single
	group multiplication:
	\begin{align*}
		CB(A_{1,2})&=CB(\sigma_1^2)\\
		&=
		\left(
		\begin{bmatrix}
			-t_1 & 1\\
			0 & 1
		\end{bmatrix},(1,2)
		\right)
		\star
		\left(
		\begin{bmatrix}
			-t_1 & 1\\
			0 & 1
		\end{bmatrix},(1,2)
		\right)\\
		&=
		\left(
		\begin{bmatrix}
			-t_1 & 1\\
			0 & 1
		\end{bmatrix}
		\cdot 
		\begin{bmatrix}
			-t_2 & 1\\
			0 & 1
		\end{bmatrix},(1,2)\cdot (1,2)
		\right)\\
		&=
		\left(
		\begin{bmatrix}
			t_1t_2 & 1-t_1\\
			0 & 1
		\end{bmatrix},e
		\right)
	\end{align*}
	and the matrix component is equal to
	$$
	\underbrace{\begin{bmatrix}
			1 & 0 \\
			0 & 1
	\end{bmatrix}}_{I_2}
	+
	\underbrace{\begin{bmatrix}
			t_2-1 & 0 \\
			0 & 0
	\end{bmatrix}}_{c_1((t_2-1)_{1\to 2})}
	+
	\underbrace{\begin{bmatrix}
			t_1t_2-t_2 & 0\\
			0 & 0
	\end{bmatrix}}_{c_1((t_1t_2-t_2)_{1\to 2})}
	+
	\underbrace{\begin{bmatrix}
			0 & 1-t_1\\
			0 & 0
	\end{bmatrix}}_{c_2((1-t_1)_{1\to 2})}
	=
	\begin{bmatrix}
		t_1t_2 & 1-t_1\\
		0 & 1
	\end{bmatrix}
	$$
	as claimed.
	\\
	\\
	Now suppose the result holds for some $B_k$ with $k>2$ and for all 
	$1\leq i<j \leq k$. Set $n=k+1$. By our inductive hypothesis, it suffices to prove 
	the formula for $j=n$. Thus, we further induct on $\ell=n-i$. When $\ell=1$, we 
	obtain the result immediately by one group multiplication:
	\begin{align*}
		CB(A_{n-1,n})&=CB(\sigma_{n-1}^2)=CB(\sigma_{n-1})\star CB(\sigma_{n-1})\\
		&=
		\left(
		\begin{bmatrix}
			I_{n-3}&\\
			& 1 & 0 & 0\\
			& t_{n-1} & -t_{n-1} & 1\\
			& 0& 0& 1
		\end{bmatrix},(n-1,n)
		\right)
		\star
		\left(
		\begin{bmatrix}
			I_{n-3}&\\
			& 1 & 0 &0\\
			& t_{n-1} & -t_{n-1}&1\\
			& 0& 0& 1
		\end{bmatrix},(n-1,n)
		\right)\\
		&=
		\left(
		\begin{bmatrix}
			I_{n-3}&\\
			& 1 & 0&0\\
			& t_{n-1} & -t_{n-1}&1\\
			& 0& 0& 1
		\end{bmatrix}
		\cdot
		\begin{bmatrix}
			I_{n-3}&\\
			& 1 & 0&0\\
			& t_n & -t_n&1\\
			& 0& 0& 1
		\end{bmatrix},(n-1,n)\cdot(n-1,n)
		\right)\\
		&=
		\left(
		\begin{bmatrix}
			I_{n-3}&\\
			& 1 & 0&0\\
			& t_{n-1}-t_{n-1}t_n & t_{n-1}t_n&1-t_{n-1}\\
			& 0& 0& 1
		\end{bmatrix}, e
		\right)
	\end{align*}
	and the matrix component is equal
	\begin{align*}
		&\underbrace{\begin{bmatrix}
				I_{n-3} & 0 & 0 & 0\\
				0 & 1 & 0 & 0 \\
				0 & 0 & 1 & 0 \\
				0 & 0 & 0 & 1
		\end{bmatrix}}_{I_n}
		+\underbrace{\begin{bmatrix}
				0 & 0 & 0 & 0\\
				0 & 0 & 0 & 0 \\
				0 & -t_{n-1}t_n+t_{n-1} &0 & 0 \\
				0 & 0 & 0 & 0
		\end{bmatrix}}_{c_{n-2}((-t_{n-1}t_n+t_{n-1})_{n-1\to n})}
		+\underbrace{\begin{bmatrix}
				0 & 0 & 0 & 0\\
				0 & 0 & 0 & 0 \\
				0 & 0 & t_n-1 & 0 \\
				0 & 0 & 0 & 0
		\end{bmatrix}}_{c_{n-1}((t_n-1)_{n-1\to n})}\\
		&+\underbrace{\begin{bmatrix}
				0 & 0 & 0 & 0\\
				0 & 0 & 0 & 0 \\
				0 & 0 & t_nt_{n-1}-t_n & 0 \\
				0 & 0 & 0 & 0
		\end{bmatrix}}_{c_{n-1}((t_nt_{n-1}-t_n)_{n-1\to n})}
		+\underbrace{\begin{bmatrix}
				0 & 0 & 0 & 0\\
				0 & 0 & 0 & 0 \\
				0 & 0 & 0 & -t_{n-1}+1 \\
				0 & 0 & 0 & 0
		\end{bmatrix}}_{c_n((-t_{n-1}+1)_{n-1\to n})}\\
		&=
		\begin{bmatrix}
			I_{n-3}&\\
			& 1 & 0&0\\
			& t_{n-1}-t_{n-1}t_n & t_{n-1}t_n&1-t_{n-1}\\
			& 0& 0& 1
		\end{bmatrix}
	\end{align*}
	as claimed.
	\\
	\\
	Now suppose our formula holds for some $2\leq \ell \leq n-2$ 
	(we treat the case $i=1$ separately). Since 
	$$A_{i-1,n}=\sigma_{i-1}A_{i,n}\sigma_{i-1}^{-1},$$ 
	we perform two group multiplications.
	Firstly, since $A_{i,n}$ is a pure braid, there is no permuting of the variables in 
	$cb(\sigma_{i-1}^{-1})$. Thus, we multiply
	\begin{align*}
		cb(A_{i,n})cb(\sigma_{i-1}^{-1})&=
		\begin{bmatrix}
			I_{i-2}&\\
			&1 & 0 & 0 & 0&0\\
			&t_i(1-t_n)&t_n & 0 & t_it_n-t_n&1-t_i\\
			&t_i(1-t_n)&t_n-1 & I_{\ell-2} & t_it_n-t_n&1-t_i\\
			&\vdots &\vdots & & \vdots&\vdots \\
			&t_i(1-t_n)&t_n-1 & 0 & 1+t_it_n-t_n &1-t_i\\
			& 0& 0& 0& 0& 1	
		\end{bmatrix}\cdot
		\begin{bmatrix}
			I_{i-3} &  &  &  \\
			& 1 & 0 & 0\\
			& 1& \frac{-1}{t_i} & \frac{1}{t_i}\\
			& 0 & 0 & 1 \\
			& & & & I_{n-i}
		\end{bmatrix}
	\end{align*}
	The product is non-unital in precisely the columns $i-1,i,i+1,$ and $n-1$ with the entries being
	$$\begin{bmatrix}
		I_{i-3}& \\
		&1 & 0 & 0  & 0 & 0 & 0\\
		&1& \frac{-1}{t_i} & \frac{1}{t_i} & 0 & 0 &0\\\\
		&t_i(1-t_n)&t_n-1 & 1 & 0 &t_it_n-t_n&1-t_i\\
		&\vdots & \vdots & I_{\ell-2}&\vdots & \vdots &\vdots \\
		&t_i(1-t_n)&t_n-1 & 0 & 0 & 1+t_it_n-t_n &1-t_i\\
		&0& 0& 0& 0& 0& 1
	\end{bmatrix}.$$
	The second group multiplication has the effect of permuting every occurrence of 
	$t_i$ in the above to $t_{i-1}$ and subsequently multiplying by $cb(\sigma_{i-1})$, 
	hence we obtain
	$${\small\begin{bmatrix}
			I_{i-3}\\
			&1& 0 & 0 \\
			&t_i & -t_i & 1\\
			&0 & 0 & 1\\
			& & & & I_{n-i}
	\end{bmatrix}}\cdot
	\begin{bmatrix}
		I_{i-3}& \\
		&1 & 0 & 0  & 0 & 0 & 0\\
		&1& \frac{-1}{t_{i-1}} & \frac{1}{t_{i-1}} & 0 & 0 &0\\\\
		&t_{i-1}(1-t_n)&t_n-1 & 1 & 0 &t_{i-1}t_n-t_n&1-t_{i-1}\\
		&\vdots & \vdots & I_{\ell-2}& \vdots & \vdots &\vdots \\
		&t_{i-1}(1-t_n)&t_n-1 & 0 & 0 & 1+t_{i-1}t_n-t_n &1-t_{i-1}\\
		&0& 0& 0& 0& 0& 1
	\end{bmatrix},$$
	which is non-unital in precisely columns $i-1,i,n-1,$ and $n$ with the entries being
	$$\begin{bmatrix}
		I_{i-2}&\\
		&1 & 0 & 0 & 0&0\\
		&t_{i-1}(1-t_n)&t_n & 0 &t_{i-1}t_n-t_n&1-t_{i-1}\\
		&t_{i-1}(1-t_n)&t_n-1 & I_{\ell-2} & t_{i-1}t_n-t_n&1-t_{i-1}\\
		&\vdots &\vdots & & \vdots &\vdots \\
		&t_{i-1}(1-t_n)&t_n-1 & 0 & 1+t_{i-1}t_n-t_n & 1-t_{i-1}\\
		&0 & 0 & 0 & 0 & 1
	\end{bmatrix}$$
	and this agrees with the proposed formula, as the reader can verify.
	\\
	\\
	Finally, we handle the case $i=1$. By induction, we have 
	$$cb(A_{2,n})=
	\begin{bmatrix}
		1 & 0 & 0 & 0&0\\
		t_2(1-t_n)&t_n & 0 &t_2t_n-t_n&1-t_2\\
		t_2(1-t_n)&t_n-1 & I_{n-4} & t_2t_n-t_n&1-t_2\\
		\vdots &\vdots & & \vdots&\vdots \\
		t_2(1-t_n)&t_n-1 & 0 & 1+t_2t_n-t_n&1-t_2\\
		0 & 0 & 0 & 0& 1
	\end{bmatrix}
	$$
	and, since $A_{1,n}=\sigma_1A_{2,n}\sigma_1^{-1}$, we obtain 
	$$cb(A_{1,n})=
	\begin{bmatrix}
		t_n & 0 &t_1t_n-t_n&1-t_1\\
		t_n-1 & I_{n-3} & t_1t_n-t_n&1-t_1\\
		\vdots & & \vdots & \vdots\\
		t_n-1 & 0 & 1+t_1t_n-t_n&1-t_1\\
		0 & 0 & 0 & 1
	\end{bmatrix},
	$$
	as the reader can verify. Finally, we note that this agrees
	with the proposed formula.
\end{proof}

\end{document}
